\documentclass[11pt]{article}
\usepackage[margin=1in]{geometry}
\usepackage{amsmath, amssymb, amsthm, bbm}
\usepackage{natbib}
\usepackage{hyperref}

\newtheorem{proposition}{Proposition}
\newcommand{\indep}{\perp\!\!\!\perp}
\newcommand{\Qmin}{Q_{\min}}
\newcommand{\Qmax}{Q_{\max}}

\title{Generating covariate-dependent left truncation times\\
under reverse-time proportional hazards}
\date{}

\begin{document}
\maketitle

This note explains how the left truncation time $Q$ is generated in the simulations of
\texttt{causalDR} (\texttt{sim\_data\_trunc()}) and of \texttt{truncAC} (the function
\texttt{gen()} in \texttt{src/gen2.R}, used for the ATE simulation of
\citealt{wang2025liberatingframeworktruncationcensoring}, hereafter Wang et al.). Both
are special cases of one construction, indexed by a baseline distribution function $G_0$
supported on an interval $[\Qmin, \Qmax]$. We explain why the construction yields a
proportional hazards model on the reversed time scale, what $\Qmin$ and $\Qmax$ control,
and how they interact with positivity and with the estimation of the truncation
distribution.

Throughout, $X = (A, Z^\top)^\top$ denotes the treatment and the covariates, $T$ is the
event time subject to left truncation (for semi-competing risks $T = \min(T_1, T_2)$), and
a subject is observed only if $Q < T$. Let $G(q \mid x) = P(Q \le q \mid X = x)$ denote the
conditional distribution function of $Q$ in the population before truncation, and let
$\eta^\top x$ be a linear predictor, possibly including an intercept $\eta_0$.

\section{Reverse-time hazards and the reverse proportional\texorpdfstring{\\}{} hazards model}

The reverse time hazard of $Q$ given $X = x$ is
\begin{equation*}
\bar\lambda(q \mid x) = \lim_{h \to 0^+}
\frac{P(q - h < Q \le q \mid Q \le q, X = x)}{h}
= \frac{g(q \mid x)}{G(q \mid x)},
\end{equation*}
where $g$ is the density of $Q$, and the reverse cumulative hazard is
$\bar H(q \mid x) = \int_q^\infty \bar\lambda(s \mid x)\,ds$. For continuous $G$, since
$\bar\lambda(s \mid x)\,ds = d\log G(s \mid x)$ and $G(\infty \mid x) = 1$, we have
$\bar H(q \mid x) = -\log G(q \mid x)$, i.e., $G(q \mid x) = \exp\{-\bar H(q \mid x)\}$,
which mirrors the relation $S = \exp(-H)$ in forward time. The reverse proportional
hazards model assumes $\bar\lambda(q \mid x) = \bar\lambda_0(q)\exp(\eta^\top x)$, which is
equivalent to $\bar H(q \mid x) = \bar H_0(q)\exp(\eta^\top x)$, and hence to
\begin{equation}\label{eq:reversePH}
G(q \mid x) = G_0(q)^{\exp(\eta^\top x)}, \qquad G_0(q) = \exp\{-\bar H_0(q)\},
\end{equation}
where $G_0$ is a baseline distribution function. Model \eqref{eq:reversePH} is the
definition of reverse proportional hazards that we use below; it remains meaningful when
$G_0$ has an atom.

The reverse proportional hazards model is a Cox model in forward time after reversing the
time axis. For any constant $\tau$ at or above the upper end of the support of $Q$, let
$W = \tau - Q$. Then $P(W > w \mid x) = P(Q < \tau - w \mid x)$, and the hazard of $W$ at
$w$ equals the reverse time hazard of $Q$ at $\tau - w$. Therefore $W$ given $X$ follows a
Cox model with baseline hazard $\bar\lambda_0(\tau - w)$ and regression coefficient $\eta$,
and this holds for every choice of $\tau$; the constant $\tau$ only fixes the origin of
the reversed time scale. Moreover, in the observed data $Q$ is right truncated by $T$, so
$W$ is left truncated by $\tau - T$: a subject is at risk at $w$ if
$\tau - T < w \le \tau - Q$, i.e., $Q \le \tau - w < T$, which is the risk set of the
backwards martingale $\bar M_Q$. This is why both packages estimate $G$ by fitting a Cox
model to the data $(\tau - T, \tau - Q]$, with $T$ replaced by $X_1$ among the subjects
with $T$ observed, and inverse probability of censoring weights.

\section{A general inverse distribution function construction}

Let $G_0$ be any distribution function supported on $[\Qmin, \Qmax]$, and let
$G_0^{-1}(v) = \inf\{q : G_0(q) \ge v\}$ be its generalized inverse, so that
$G_0^{-1}(v) \le q$ if and only if $v \le G_0(q)$.

\begin{proposition}\label{prop:construction}
Let $U \sim \mathrm{Uniform}(0, 1)$ be independent of $X$ and of the potential event times,
let $V = U^{\exp(-\eta^\top X)}$, and let $Q = G_0^{-1}(V)$. Then $Q$ given $X$ follows the
reverse proportional hazards model \eqref{eq:reversePH}, and $Q$ is independent of the
potential event times given $X$.
\end{proposition}

\begin{proof}
For $q \in [\Qmin, \Qmax]$,
\begin{equation*}
P(Q \le q \mid X = x) = P\{V \le G_0(q) \mid X = x\}
= P\{U \le G_0(q)^{\exp(\eta^\top x)}\} = G_0(q)^{\exp(\eta^\top x)},
\end{equation*}
where the second equality raises both sides of $U^{\exp(-\eta^\top x)} \le G_0(q)$ to the
power $\exp(\eta^\top x) > 0$. Conditional independence holds because $Q$ is a function of
$X$ and of $U$, and $U$ is generated independently of everything else.
\end{proof}

In words, $V = G(Q \mid X)$ is the conditional probability integral transform of $Q$: we
first draw $G(Q \mid X)$, which is uniform, and then invert $G(\cdot \mid X)$. This is the
line \texttt{G0\_Q = runif(N)\^{}exp(-lp)} in \texttt{sim\_data\_trunc()}. The conditional
independence of $Q$ and the event times given $X$ is the conditional independent
truncation assumption.

Two remarks on $\eta$. First, in the Cox model for $W = \tau - Q$ an intercept $\eta_0$ is
not identifiable, since $G_0^{\exp(\eta_0 + \eta_1^\top x)} = \{G_0^{\exp(\eta_0)}\}^{
\exp(\eta_1^\top x)}$ and $G_0^{\exp(\eta_0)}$ is another baseline distribution function;
in the simulation, however, $\eta_0$ changes the distribution of $Q$ and is a convenient
way to tune the truncation rate. Second, a larger $\eta^\top x$ makes the exponent
$\exp(-\eta^\top x)$ of $U$ smaller, so $V$ is closer to 1 and $Q$ is stochastically larger;
subjects with a larger linear predictor are therefore more likely to be truncated.

\section{The uniform baseline used by Wang et al.}

Wang et al.\ take the baseline to be the uniform distribution on $[\Qmin, \Qmax]$,
\begin{equation*}
G_0(q) = \frac{q - \Qmin}{\Qmax - \Qmin}, \qquad \Qmin \le q \le \Qmax,
\end{equation*}
so that by Proposition \ref{prop:construction}
\begin{equation}\label{eq:uniform}
Q = \Qmin + (\Qmax - \Qmin)\,U^{\exp(-\eta^\top X)}, \qquad
G(q \mid x) = \left\{\frac{q - \Qmin}{\Qmax - \Qmin}\right\}^{\exp(\eta^\top x)}.
\end{equation}
Their code generates $Q$ through the reversed time $Q_2 = \tau - Q$ instead:
\begin{equation*}
Q_2 = Q_{2,\min} + (Q_{2,\max} - Q_{2,\min})\{1 - U^{\exp(-\eta^\top X)}\},
\qquad Q = \tau - Q_2,
\end{equation*}
with $Q_{2,\max} = \tau - \Qmin$ and $Q_{2,\min} = \tau - \Qmax$. Substituting these
expressions gives $Q = \Qmax - (\Qmax - \Qmin)\{1 - U^{\exp(-\eta^\top X)}\}$, which is
\eqref{eq:uniform}. Hence $\tau$ cancels and plays no role in the distribution of $Q$.
The description ``$\tau - Q$ follows a proportional hazards model with the baseline hazard
of a uniform distribution'' follows from Section 1: when $\eta^\top x = 0$, $W = \tau - Q$
is uniform on $[\tau - \Qmax, \tau - \Qmin]$, whose hazard is
$\lambda_0(w) = 1/(\tau - \Qmin - w)$, and the hazard of $W$ given $X = x$ is
$\lambda_0(w)\exp(\eta^\top x)$.

The parameters $\Qmin$ and $\Qmax$ are the lower and upper ends of the support of $Q$:
$G(q \mid x) = 0$ for $q < \Qmin$ and $G(q \mid x) = 1$ for $q \ge \Qmax$, for every $x$.
They allow the truncation window to be moved, e.g., $\Qmin > 0$ when all subjects enter the
study late, or narrowed, e.g., $\Qmax$ small relative to the range of $T$ for light
truncation. Since $E(U^c) = 1/(1 + c)$ for $c > 0$,
\begin{equation*}
E(Q \mid X = x) = \Qmin + \frac{\Qmax - \Qmin}{1 + \exp(-\eta^\top x)}.
\end{equation*}
In the ATE simulation of Wang et al., $\Qmin = 0$, $\Qmax = \tau = 5$, and
$\eta^\top x = -0.6 + 0.6a + 0.4z_1 + 0.2z_2$, so that, e.g.,
$E(Q \mid A = 0, Z = 0) = 5/\{1 + \exp(0.6)\} = 1.77$ and $E(Q \mid A = 1, Z = 0) = 2.5$.
With the uniform baseline $Q$ is continuous, and there is no probability mass at $\Qmin$.

\section{A baseline with a point mass at the lower end}

The function \texttt{sim\_data\_trunc()} in \texttt{causalDR} uses the baseline
\begin{equation*}
G_0(q) = p_0 + (1 - p_0)\frac{q - \Qmin}{\Qmax - \Qmin}, \qquad \Qmin \le q \le \Qmax,
\end{equation*}
with $\Qmin = 0$ and $\Qmax$ given by the argument \texttt{q\_max}. The generalized inverse
is $G_0^{-1}(v) = \Qmin$ for $v \le p_0$ and
$G_0^{-1}(v) = \Qmin + (\Qmax - \Qmin)(v - p_0)/(1 - p_0)$ for $v > p_0$, so that, by
Proposition \ref{prop:construction}, $Q$ follows the reverse proportional hazards model
with
\begin{equation*}
P(Q = \Qmin \mid X = x) = G(\Qmin \mid x) = p_0^{\exp(\eta^\top x)}.
\end{equation*}
The uniform baseline of Wang et al.\ is the special case $p_0 = 0$. The point mass
represents subjects who are followed from the time origin, i.e., incident rather than
prevalent subjects; such subjects are never truncated, since $\Qmin < T$.

Other familiar constructions can be checked in the same way. If
$W \sim \mathrm{Exp}\{\lambda\exp(\eta^\top x)\}$ given $X = x$ and
$Q = \max(\tau - W, 0)$, then for $0 \le q \le \tau$
\begin{equation*}
G(q \mid x) = P(W \ge \tau - q \mid x) = \exp\{-\lambda e^{\eta^\top x}(\tau - q)\}
= G_0(q)^{\exp(\eta^\top x)}, \qquad G_0(q) = \exp\{-\lambda(\tau - q)\},
\end{equation*}
which is again of the form \eqref{eq:reversePH}, with a point mass
$\exp\{-\lambda\tau e^{\eta^\top x}\}$ at 0. In contrast, generating $W$ from the same
exponential distribution conditional on $W \le \tau$, or generating $Q$ itself from an
exponential distribution with rate $\exp(\eta^\top x)$, does not give a reverse
proportional hazards model, because
$-\log G(q \mid x)$ is then not of the form $\bar H_0(q)\exp(\eta^\top x)$; these
constructions can be used to misspecify the model for $G$.

A baseline $G_0$ without an atom at $\Qmin$ requires $G_0(\Qmin) = 0$, i.e., the reverse
cumulative hazard $\bar H_0$ must diverge at $\Qmin$, so that $W = \tau - Q$ has a bounded
support ending at $\tau - \Qmin$ with a hazard that diverges there. The uniform baseline
has this property, whereas the exponential construction above, whose support is
unbounded before capping, cannot avoid the atom.

\section{Forward-time proportional hazards for the event times}

The same idea generates event times from a proportional hazards model in forward time,
\begin{equation}\label{eq:forwardPH}
S(t \mid x) = P(T > t \mid X = x) = S_0(t)^{\exp(\eta^\top x)}, \qquad
H(t \mid x) = H_0(t)\exp(\eta^\top x),
\end{equation}
where $S_0 = \exp(-H_0)$ is a baseline survival function, with the survival function in
place of the distribution function. Let $S_0^{-1}(v) = \inf\{t : S_0(t) \le v\}$ be the
generalized inverse of $S_0$, so that, since $S_0$ is non-increasing and right-continuous,
$S_0^{-1}(v) \le t$ if and only if $S_0(t) \le v$.

\begin{proposition}\label{prop:forward}
Let $U \sim \mathrm{Uniform}(0, 1)$ be independent of $X$, and let
$T = S_0^{-1}\{U^{\exp(-\eta^\top X)}\}$. Then $T$ given $X$ follows the proportional
hazards model \eqref{eq:forwardPH}.
\end{proposition}

\begin{proof}
For $t \ge 0$,
\begin{equation*}
P(T \le t \mid X = x) = P\{U^{\exp(-\eta^\top x)} \ge S_0(t)\}
= P\{U \ge S_0(t)^{\exp(\eta^\top x)}\} = 1 - S_0(t)^{\exp(\eta^\top x)},
\end{equation*}
so that $S(t \mid x) = S_0(t)^{\exp(\eta^\top x)}$.
\end{proof}

In words, we first draw $S(T \mid X)$, which is uniform, and then invert
$S(\cdot \mid X)$. Since $E = -\log U$ follows the standard exponential distribution, the
construction can equivalently be written as $T = H_0^{-1}\{E\exp(-\eta^\top X)\}$, i.e., the
cumulative hazard $H(T \mid X)$ is standard exponential, which is the method of
\citet{bender2005generating} for simulating Cox models. For example, the exponential
baseline $H_0(t) = \lambda t$ gives $T = -\log(U)\exp(-\eta^\top X)/\lambda$, the Weibull
baseline $H_0(t) = \lambda t^\nu$ gives
$T = \{-\log(U)\exp(-\eta^\top X)/\lambda\}^{1/\nu}$, and the uniform baseline
$S_0(t) = 1 - t/\tau$ on $[0, \tau]$ gives $T = \tau\{1 - U^{\exp(-\eta^\top X)}\}$. In
contrast to the reverse-time construction, a larger $\eta^\top x$ now makes
$U^{\exp(-\eta^\top x)}$ closer to 1 and hence $T$ stochastically smaller, since a larger
hazard leads to earlier events. A location shift $T_{\min} + S_0^{-1}\{U^{\exp(-\eta^\top X)}\}$
keeps the proportional hazards structure, with all hazards equal to zero on
$[0, T_{\min})$.

The two constructions are mirror images of each other. By Section 1, $Q$ follows the
reverse proportional hazards model \eqref{eq:reversePH} if and only if $W = \tau - Q$
follows the forward model \eqref{eq:forwardPH} with the baseline survival function
$S_0(w) = G_0\{(\tau - w)-\}$, and for the uniform baseline the forward draw
$W = \tau\{1 - U^{\exp(-\eta^\top X)}\}$ is exactly the reversed-time draw
$Q_2 = \tau - Q$ in the code of Wang et al.\ with $\Qmin = 0$ and $\Qmax = \tau$. In the
semi-competing risks generator of \texttt{sim\_data\_trunc()}, the latent event time
$T_{\min} + \Lambda_0^{a,-1}\{-\log(U_1)/2\}$ is of this form with $H_0 = \Lambda_0^a$ and the
constant multiplier $\exp(\eta^\top x) = 2$; the covariates enter only through their
correlation with the latent $(U_1, U_2)$, which is why the Cox models for the transition
rates given $(A, Z)$ are misspecified in that design.

\section{Positivity and the choice of \texorpdfstring{$\Qmin$}{Q\_min} and \texorpdfstring{$\Qmax$}{Q\_max}}

The estimators weight the observed subjects by $1/G(T \mid A, Z)$, so $G(T \mid A, Z)$ must
be bounded away from zero; Wang et al.\ assume that $T \ge \tau_1$ and $Q < \tau_2$ almost
surely, $G(\tau_1 \mid A, Z) \ge \epsilon$ and $1 - F(\tau_2 \mid A, Z) \ge \epsilon$ for
some $\epsilon > 0$, where $F$ is the distribution function of $T$ given $(A, Z)$.

The lower end $\Qmin$ must lie below the lower end of the support of $T$. Under the
uniform baseline and $T \ge T_{\min}$,
\begin{equation*}
G(T \mid x) \ge G(T_{\min} \mid x)
= \left\{\frac{T_{\min} - \Qmin}{\Qmax - \Qmin}\right\}^{\exp(\eta^\top x)},
\end{equation*}
which is bounded away from zero if and only if $T_{\min} > \Qmin$ and $\eta^\top x$ is
bounded. If instead $T$ can be arbitrarily close to $\Qmin$, e.g., $T_{\min} = \Qmin = 0$,
then $G(T \mid x) \to 0$ as $T \to \Qmin$ and the weights are unbounded. There are two ways
to secure positivity. Wang et al.\ shift the event time, $T = 1 + \text{Weibull}$, so that
$T_{\min} = \tau_1 = 1 > \Qmin = 0$ and $G(T \mid x) \ge 0.2^{\exp(\eta^\top x)}$. The
point mass of Section 4 achieves the same without shifting $T$, since
$G(T \mid x) \ge G(\Qmin \mid x) = p_0^{\exp(\eta^\top x)}$. For semi-competing risks,
shifting $T = \min(T_1, T_2)$ amounts to shifting both $T_1$ and $T_2$ by $\tau_1$, i.e.,
all transition hazards vanish on $[0, \tau_1)$; Cox models for the transition rates remain
correctly specified, since their baseline hazards are simply zero on that interval.

The upper end $\Qmax$ plays the role of $\tau_2$. A subject with $Q$ close to $\Qmax$ is
observed only if $T > Q$, so $G$ is identified on its whole support only if
$P(T > \Qmax \mid x) > 0$. The upper end also controls the amount of truncation: by the
conditional independence of $Q$ and $T$,
\begin{equation*}
P(Q < T \mid X = x) = E\{G(T- \mid x) \mid X = x\},
\end{equation*}
which decreases as $\Qmax$ increases with $\Qmin$ fixed, because $G_0(q)$ then decreases
for every $q$, and as $\eta^\top x$ increases. In the ATE simulation of Wang et al.\ about
23\% of the generated subjects are truncated.

\section{Implications for the estimation of \texorpdfstring{$G$}{G}}

With a continuous baseline such as the uniform one there are no ties among the
reversed-time event times $\tau - Q$, and the Cox model for $W = \tau - Q$ described in
Section 1 is fitted without complications.

With a point mass at $\Qmin$, all subjects with $Q = \Qmin$ have their reversed-time event
at $w = \tau - \Qmin$, and the risk set at that time, $\{Q \le \Qmin < T\}$, consists of
exactly these subjects. The true reverse time hazard at $\Qmin$ is
$P(Q = \Qmin \mid Q \le \Qmin, X = x) = 1$ for every $x$, so the exact partial likelihood
contribution of this tied event time is 1 and carries no information about $\eta$.
The Breslow approximation replaces it by
\begin{equation*}
\frac{\prod_{i \in \mathcal{D}} \exp(\eta^\top x_i)}
{\left\{\sum_{j \in \mathcal{D}} \exp(\eta^\top x_j)\right\}^{|\mathcal{D}|}},
\end{equation*}
where $\mathcal{D}$ is the set of subjects with $Q = \Qmin$, and this factor is maximized
at $\eta = 0$ (the Efron approximation has the same dependence on $\eta$); it acts as a
penalty that shrinks the estimate of $\eta$ towards zero and does not vanish as the sample
size grows. The information carried by $Q = \Qmin$ is that no reversed-time event occurs
before $w = \tau - \Qmin$, since $G(\Qmin \mid x) = \exp\{-e^{\eta^\top x}\bar H_0(\Qmin)\}$
and $\bar H_0(\Qmin)$ only involves the reverse time hazard above $\Qmin$. The fit should
therefore treat the subjects with $Q = \Qmin$ as censored at the end of the reversed time
scale, i.e., use the event indicator $\mathbbm{1}(Q > \Qmin)$, and set
$G(\Qmin- \mid x) = 0$, so that the atom $dG(\Qmin \mid x) = G(\Qmin \mid x)$ is retained
in the integrals with respect to the backwards martingale $\bar M_Q$.

\bibliographystyle{apalike}
\bibliography{truncation_time_generation}

\end{document}
